% Abstract

日本のアニメ制作市場は2024年に過去最高の3{,}621億円（24.1億ドル）に達したが、制作労働力の47--70\%を占めるフリーランスアニメーターは、下請け構造の階層化、スタッフの流動性、クレジット枠の制約により、公式クレジットにおいて依然として大部分が不可視のままである。本論文は、Weyl et al.~(2022)が提唱したソウルバウンドトークン（Soulbound Token, SBT）---特定のブロックチェーンアドレスに紐づけられた譲渡不可能なトークン---を活用し、個々の制作タスク単位でアニメーターの貢献を不変かつ検証可能な形で記録するブロックチェーンベースの貢献証明フレームワーク\textbf{AnimatorSBT}を提案する。本フレームワークは、フリーランスアニメーター向けの既存の制作管理ツールとの統合を前提に設計されており、Polygon PoSブロックチェーン上への展開を対象としている。パブリックテストネットにデプロイされたスマートコントラクト、発行サービス、および検証ポータルから構成されるプロトタイプにより設計を検証し、ガスコストを実測した。実測値を用いたコストシミュレーションと要件分析により、本システムが経済的に実現可能であること（実測された発行コストは証明書1件あたり\$0.001未満）、および理想的な資格証明ソリューションの5つの要件---完全性、検証可能性、永続性、可搬性、最小摩擦---を満たすことを示す。本フレームワークは、クリエイティブ分野のギグワーカーに対する分散型アイデンティティ（Decentralized Identity）に関する新興の研究に貢献するとともに、他のプロジェクトベースのクリエイティブ産業に適用可能なモデルを提示するものである。
